

\lecture{2}

\lecturedate{22 января 2026}

\section{Предел последовательности}

\begin{definition}
Число $A$ называется пределом последовательности $\{x_n\}$, если для любого $\varepsilon > 0$ существует номер $N(\varepsilon)$, такой что для всех $n > N$ выполняется неравенство:
\[
|x_n - A| < \varepsilon
\]
\end{definition}

Обозначение: $\lim_{n \to \infty} x_n = A$.

\subsection{Свойства сходящихся последовательностей}
\begin{theorem}
Сходящаяся последовательность имеет только один предел.
\end{theorem}

\begin{proof}
Предположим противное, что существуют два предела $A$ и $B$ ($A \neq B$). Выберем $\varepsilon = \frac{|A-B|}{2}$. Тогда окрестности точек не пересекаются, что приводит к противоречию с определением предела.
\end{proof}


